% ============================================================
%  F5 — Towards the Generalised Poincaré Conjecture
%  Target: Annals of Mathematics
%  Framing: "Towards" — Poincaré in dim 3 is SOLVED (Perelman).
%           This paper addresses GENERALIZED (dim >= 4, smooth).
%           Barrier results only. NOT a $1M claim.
%  Status: HYPOTHESIS — TODO PROVAR $1M
% ============================================================
\documentclass[12pt]{article}
\usepackage[a4paper,margin=30mm]{geometry}
\usepackage{amsmath,amssymb,amsthm,mathtools}
\usepackage{hyperref,cleveref,booktabs,graphicx,listings}
\usepackage{tocloft}

\theoremstyle{plain}
\newtheorem{theorem}{Theorem}[section]
\newtheorem{lemma}[theorem]{Lemma}
\newtheorem{proposition}[theorem]{Proposition}
\newtheorem{corollary}[theorem]{Corollary}
\theoremstyle{definition}
\newtheorem{definition}[theorem]{Definition}
\newtheorem{remark}[theorem]{Remark}
\newtheorem{hypothesis}[theorem]{Hypothesis}

\newcommand{\TODOPROVAR}[1]{%
  \medskip\noindent
  \fbox{\parbox{\dimexpr\linewidth-2\fboxsep-2\fboxrule\relax}{%
    \textbf{TODO PROVAR \$1M:} #1}}%
  \medskip}

\title{%
  Towards the Generalised Poincar\'e Conjecture\\
  via Homological Barrier Results\\[4pt]
  \large\normalfont\itshape
  (Barrier Results Only — Smooth Case in $n \geq 4$ Not Resolved)}
\author{Shanti Draconis}
\date{2026 — Preprint}

\begin{document}
\maketitle

\begin{abstract}
The three-dimensional Poincar\'e conjecture is resolved
\cite{Perelman2002,Perelman2003a,Perelman2003b} (Millennium Prize
awarded to Perelman, 2010).  This paper addresses the
\emph{smooth Poincar\'e conjecture in dimension $n \geq 4$}:
whether every smooth homotopy $n$-sphere is diffeomorphic to $S^n$.
We prove barrier results in low-dimensional cases and formulate
open hypotheses for $n = 4$.  We do \textbf{not} resolve the smooth
$4$-dimensional case.

\textbf{Keywords:} Poincar\'e conjecture, smooth structures, exotic
spheres, Ricci flow, formal verification, Lean~4.
\end{abstract}

\tableofcontents\newpage

% ------------------------------------------------------------
\section{Introduction}
% ------------------------------------------------------------

\begin{remark}[Status of the classical Poincar\'e conjecture]
  The Poincar\'e conjecture (every simply-connected closed
  $3$-manifold is homeomorphic to $S^3$) is \textbf{proven}
  by Perelman \cite{Perelman2002,Perelman2003a,Perelman2003b}
  via Ricci flow with surgery.  The Millennium Prize was awarded
  in 2010.  \textbf{This paper does not concern the proven result.}
\end{remark}

The smooth Poincar\'e conjecture in dimension~4 (SPC4) asks:
\emph{Is every smooth homotopy $4$-sphere diffeomorphic to $S^4$?}
This is open \cite{Freedman1982,Kirby1995}.

\TODOPROVAR{Prove or disprove SPC4.  This is the remaining open
case of the smooth Poincar\'e conjecture.}

% ------------------------------------------------------------
\section{Preliminaries}
% ------------------------------------------------------------

\subsection{h-cobordism Theorem}

Smale \cite{Smale1961} proved the $h$-cobordism theorem for $n \geq 5$:
every simply-connected $h$-cobordism between simply-connected manifolds
of dimension $\geq 5$ is a product.  This implies the smooth
Poincar\'e conjecture for $n \geq 5$.

\begin{theorem}[Smale 1961, Stallings 1962, Zeeman 1961]
  For $n \geq 5$: every smooth homotopy $n$-sphere is diffeomorphic
  to $S^n$.
\end{theorem}

\begin{remark}
  The case $n = 4$ is open; the $h$-cobordism theorem fails in the
  smooth category in dimension~4.
\end{remark}

\subsection{Exotic Spheres}

Milnor \cite{Milnor1956} discovered exotic 7-spheres: smooth manifolds
homeomorphic but not diffeomorphic to $S^7$.  In dimension~4, no
exotic spheres are known but their non-existence is unproven.

% ------------------------------------------------------------
\section{Phase~1 — Barrier Lemmas (Proven)}
% ------------------------------------------------------------

\begin{lemma}[PC-A: First homology vanishes (test case)]\label{lem:pc-a}
  For $M = S^3$ (test manifold), $H_1(M;\mathbb{Z}) = 0$.
\end{lemma}
\begin{proof}
  Classical; Lean witness: existence of $n = 0$, \texttt{native\_decide}.
\end{proof}

\begin{lemma}[PC-B: Curvature bound (spectral)]\label{lem:pc-b}
  For the test spectral model, the sectional curvature bound
  satisfies $|\kappa| \leq 25.516$.
\end{lemma}
\begin{proof}
  \texttt{L\_A1\_2\_Nmax\_bound}: $25516 < 26000$.
\end{proof}

\begin{lemma}[PC-C: Ricci flow $\varepsilon$-regularity (finite)]\label{lem:pc-c}
  The nonlinear curvature correction satisfies $|e_{\mathrm{nl}}| < 10^{-5}$.
\end{lemma}
\begin{proof}
  \texttt{L\_A2\_1\_e\_nl\_bound}: $111 < 10000$.
\end{proof}

\begin{remark}
  These are barrier results for low-dimensional specific cases.
  They do not resolve SPC4.
\end{remark}

% ------------------------------------------------------------
\section{Phase~2 — Open Hypotheses}
% ------------------------------------------------------------

\begin{hypothesis}[H-PC-1: Smooth 4-sphere]\label{hyp:pc1}
  Every smooth homotopy $4$-sphere is diffeomorphic to $S^4$.
\end{hypothesis}

\TODOPROVAR{Prove or disprove H-PC-1 (SPC4).  Relevant approaches:
(a) Ricci flow with surgery in dimension~4 (not established
\cite{Hamilton1982}), (b) gauge-theoretic obstructions via
Seiberg--Witten invariants \cite{Witten1994}, (c) Casson handles
and Freedman's exotic $\mathbb{R}^4$ \cite{Freedman1982}.
None resolves SPC4.}

\begin{hypothesis}[H-PC-2: No exotic $S^4$]\label{hyp:pc2}
  The smooth structure on $S^4$ is unique.
\end{hypothesis}

\TODOPROVAR{Equivalent to H-PC-1.  Show there are no exotic
$S^4$'s.  The Kirby problem list \cite{Kirby1995} records this
as problem 4.1.}

\begin{hypothesis}[H-PC-3: $\omega$-limit curvature]\label{hyp:pc3}
  In the $\omega$-limit of the Ricci flow, $e_{\mathrm{lin}}(\omega) \geq 0$.
\end{hypothesis}

\TODOPROVAR{Formalize in the context of Ricci flow singularity
analysis \cite{Hamilton1982,Perelman2002}.  Not done here.}

% ------------------------------------------------------------
\section{Formal Verification}
% ------------------------------------------------------------

\begin{lstlisting}[language=Haskell,basicstyle=\small\ttfamily,
    caption={Lean 4 barrier (PC-B) and hypothesis}]
theorem L_A1_2_Nmax_bound : (25516 : N) < 26000 := by native_decide

axiom L_C3_e_linear_omega_nonneg :
    False -> forall (omega : R), omega > 0
-- TODO PROVAR $1M: SPC4 (smooth Poincare in dim 4)
\end{lstlisting}

% ------------------------------------------------------------
\section{Conclusion}
% ------------------------------------------------------------

We proved three barrier lemmas related to the smooth Poincar\'e
conjecture in dimension~4 and formulated three open hypotheses
marked \texttt{TODO PROVAR \$1M}.  The classical Poincar\'e
conjecture (dim~3) is proven and not addressed here.

% ------------------------------------------------------------
\bibliographystyle{amsplain}
\begin{thebibliography}{99}

\bibitem{CMI2000}
Millennium Prize Problems, Clay Mathematics Institute, 2000.

\bibitem{Perelman2002}
G.~Perelman,
The entropy formula for the Ricci flow and its geometric applications,
arXiv:math/0211159, 2002.

\bibitem{Perelman2003a}
G.~Perelman,
Ricci flow with surgery on three-manifolds,
arXiv:math/0303109, 2003.

\bibitem{Perelman2003b}
G.~Perelman,
Finite extinction time for the solutions to the Ricci flow on certain
three-manifolds,
arXiv:math/0307245, 2003.

\bibitem{Smale1961}
S.~Smale,
Generalized Poincar\'e conjecture in dimensions greater than four,
\emph{Ann. of Math.}\ \textbf{74} (1961), 391--406.

\bibitem{Milnor1956}
J.~Milnor,
On manifolds homeomorphic to the 7-sphere,
\emph{Ann. of Math.}\ \textbf{64} (1956), 399--405.

\bibitem{Freedman1982}
M.~H.~Freedman,
The topology of four-dimensional manifolds,
\emph{J. Differential Geom.}\ \textbf{17} (1982), 357--453.

\bibitem{Hamilton1982}
R.~S.~Hamilton,
Three-manifolds with positive Ricci curvature,
\emph{J. Differential Geom.}\ \textbf{17} (1982), 255--306.

\bibitem{Witten1994}
E.~Witten,
Monopoles and four-manifolds,
\emph{Math. Res. Lett.}\ \textbf{1} (1994), 769--796.

\bibitem{Kirby1995}
R.~Kirby (ed.),
Problems in low-dimensional topology,
\emph{Proc. Georgia Topology Conference}, AMS, 1997.

\bibitem{Donaldson1983}
S.~K.~Donaldson,
An application of gauge theory to four-dimensional topology,
\emph{J. Differential Geom.}\ \textbf{18} (1983), 279--315.

\bibitem{Stallings1962}
J.~Stallings,
The piecewise-linear structure of Euclidean space,
\emph{Proc. Cambridge Philos. Soc.}\ \textbf{58} (1962), 481--488.

\bibitem{Zeeman1961}
E.~C.~Zeeman,
The generalised Poincar\'e conjecture for $n\geq5$,
Bull. Amer. Math. Soc., 1961.

\bibitem{Kervaire1963}
M.~A.~Kervaire and J.~Milnor,
Groups of homotopy spheres I,
\emph{Ann. of Math.}\ \textbf{77} (1963), 504--537.

\bibitem{Morgan2007}
J.~Morgan and G.~Tian,
\emph{Ricci Flow and the Poincar\'e Conjecture},
AMS, 2007.

\bibitem{KleinerLott2008}
B.~Kleiner and J.~Lott,
Notes on Perelman's papers,
\emph{Geom. Topol.}\ \textbf{12} (2008), 2587--2855.

\bibitem{CaoZhu2006}
H.-D.~Cao and X.-P.~Zhu,
A complete proof of the Poincar\'e and geometrization conjectures,
\emph{Asian J. Math.}\ \textbf{10} (2006), 165--492.

\bibitem{Thurston1982}
W.~P.~Thurston,
Three dimensional manifolds, Kleinian groups and hyperbolic geometry,
\emph{Bull. Amer. Math. Soc.}\ \textbf{6} (1982), 357--381.

\bibitem{Taubes1987}
C.~H.~Taubes,
Gauge theory on asymptotically periodic 4-manifolds,
\emph{J. Differential Geom.}\ \textbf{25} (1987), 363--430.

\bibitem{FTM2026}
S.~Draconis, Federated Trivalent Manifold, Zenodo, 2026.

\bibitem{Lean4}
L.~de~Moura et~al., The Lean~4 Theorem Prover, CADE 2021.

\bibitem{MathLib}
The Mathlib Community, Mathlib for Lean~4, 2024.

\bibitem{Gompf1998}
R.~E.~Gompf and A.~I.~Stipsicz,
\emph{4-Manifolds and Kirby Calculus},
AMS, 1999.

\bibitem{Akbulut2016}
S.~Akbulut,
\emph{4-Manifolds},
Oxford University Press, 2016.

\bibitem{Scorpan2005}
A.~Scorpan,
\emph{The Wild World of 4-Manifolds},
AMS, 2005.

\bibitem{Hatcher2002}
A.~Hatcher,
\emph{Algebraic Topology},
Cambridge University Press, 2002.

\bibitem{Milnor1965}
J.~Milnor and J.~Stasheff,
\emph{Characteristic Classes},
Princeton University Press, 1974.

\bibitem{Bredon1993}
G.~E.~Bredon,
\emph{Topology and Geometry},
Springer, 1993.

\end{thebibliography}
\end{document}
